\documentclass{article}

% ICLR 2026 style - download from https://github.com/ICLR/Master-Template/raw/master/iclr2026.zip
% and place iclr2026_conference.sty in the same folder
\usepackage{iclr2026_conference, times}

% Standard packages
\usepackage{hyperref}
\usepackage{url}
\usepackage{booktabs}
\usepackage{amsmath}
\usepackage{amssymb}
\usepackage{graphicx}
\usepackage{microtype}
\usepackage{xcolor}
\usepackage{multirow}
\usepackage{array}

\title{AdaptiSelect-TSP: Phase-Aware Q-Learning for\\
Adaptive Local Search Operator Selection\\
on the Travelling Salesman Problem}

\author{
  S{\i}la Erda\c{s} \quad Özlem Akyaz \\
  Department of Industrial Engineering \\
  Hacettepe University, Ankara, Turkey \\
  \texttt{\{sila.erdas, ozlem.akyaz\}@hacettepe.edu.tr}
}

\newcommand{\fix}{\marginpar{FIX}}
\newcommand{\new}{\marginpar{NEW}}

\iclrfinalcopy % Remove this for anonymous submission

\begin{document}

\maketitle

% ─────────────────────────────────────────────────────────────────────────────
\begin{abstract}
The Travelling Salesman Problem (TSP) is a canonical NP-hard combinatorial
optimization problem. Local search operators such as 2-opt, swap, and relocate
are effective individually, yet their relative merits shift dramatically as the
search progresses: 2-opt excels at untangling early-stage tours; relocate
refines near-optimal solutions; and restart is the only reliable escape from
deep local minima. Committing to a single operator throughout the search is
therefore inherently suboptimal. We present \textbf{AdaptiSelect-TSP}, a
tabular Q-learning hyper-heuristic that learns a phase-aware operator selection
policy for the TSP. The Markov Decision Process (MDP) encodes solution quality,
iteration stage, recent improvement, and last operator into 120 discrete states.
We compare two reward shapings (pure improvement R1 vs.\ global-best bonus R2),
three epsilon management strategies including an adaptive stagnation-driven
schedule, and provide an empirical operator transition matrix for synergy
analysis. On four TSPLIB benchmark instances (eil51, berlin52, st70, kroA100),
the Q-learning controller achieves optimality gaps of 3.00\%, 1.03\%, 3.82\%,
and 3.90\%, respectively, outperforming fixed 2-opt by up to 1.6 percentage
points and random operator selection by up to 12.8 percentage points.
\end{abstract}

% ─────────────────────────────────────────────────────────────────────────────
\section{Introduction}

The Travelling Salesman Problem (TSP) underpins applications ranging from
vehicle routing and logistics to circuit board drilling and DNA sequencing.
Despite decades of research, exact solvers scale poorly beyond a few hundred
cities, making heuristic and meta-heuristic algorithms, particularly local
search procedures, indispensable in practice~\citep{reinelt1991tsplib}.
Standardized benchmarks such as TSPLIB~\citep{reinelt1991tsplib} have
established the TSP as the canonical testbed for evaluating new search
methodologies, a role it continues to play for the rapidly growing body of
machine-learning-augmented optimization~\citep{mazyavkina2021rl,chung2025neural}.

Classical local search commits to a single neighborhood operator applied
repeatedly until a local optimum is reached. Each operator has a characteristic
regime of effectiveness: \textit{2-opt} is most productive early, when tours
contain many long crossing edges; \textit{relocate} (Or-opt) excels at refining
nearly optimal tours; \textit{swap} provides fast but coarse diversification;
and \textit{restart} is sometimes the only escape from deep basins. A
fixed-operator strategy therefore performs well only in one phase of the search
trajectory. The experimental results of this paper confirm this intuition: fixed
swap and relocate produce optimality gaps of 17--28\% on benchmark instances,
while fixed 2-opt achieves only 2--5\% --- yet the Q-learning controller, by
dynamically selecting among all four, consistently surpasses fixed 2-opt.

This observation motivates \textit{selection hyper-heuristics}: high-level
controllers that learn to choose which low-level heuristic to apply at each
iteration based on contextual feedback~\citep{burke2013hyper}. Among selection
mechanisms, Q-learning~\citep{watkins1992qlearning} provides a principled
framework for learning state-conditioned action preferences from delayed
rewards. The contributions of this paper are fourfold:

\begin{enumerate}
\item \textbf{Phase-aware MDP formulation.} We design a 120-state MDP encoding
solution quality gap, iteration stage, recent improvement flag, and last
operator, enabling the policy to adapt across early, mid, and late search phases
without neural function approximation.

\item \textbf{Reward shaping ablation.} We empirically compare pure improvement
reward (R1) with a global-best bonus reward (R2), isolating the contribution of
bonus shaping to convergence speed and final solution quality.

\item \textbf{Adaptive $\varepsilon$-greedy with stagnation detection.} We
propose and benchmark a stagnation-driven exploration schedule that increases
$\varepsilon$ after a threshold of non-improving steps and contracts it upon
improvement, compared against constant and exponentially decaying baselines.

\item \textbf{Operator synergy analysis.} We introduce an empirical operator
transition matrix recording conditional success frequencies, providing
quantitative diagnostics for operator complementarity.
\end{enumerate}

% ─────────────────────────────────────────────────────────────────────────────
\section{Related Work}

\subsection{TSP and Local Search}

Local search for the TSP has been studied since the introduction of 2-opt and
3-opt edge-exchange neighborhoods and the Lin--Kernighan (LK) variable-depth
heuristic~\citep{lin1973lk}, which for decades represented the state of the art
in symmetric TSP heuristics. \citet{reinelt1991tsplib} established the TSPLIB
benchmark library, the standard test bed for TSP heuristics used in this paper.
Subsequent work confirmed that combinations of operators --- for example 2-opt
with relocate under iterated local search --- substantially outperform any
single operator, motivating adaptive selection mechanisms.

\subsection{Hyper-Heuristics}

The hyper-heuristic paradigm was introduced by \citet{cowling2001hyper} with a
choice-function controller balancing reward, frequency, and recency.
\citet{burke2013hyper} provide a comprehensive survey classifying
hyper-heuristics along selection vs.\ generation and constructive vs.\
perturbative dimensions. Within selection hyper-heuristics, Q-learning has been
a recurring choice: \citet{karimim2021ils} integrate Q-learning into Iterated
Local Search for TSP, demonstrating that intelligently selecting operators yields
lower optimality gaps and faster convergence than fixed strategies. A general
lesson from this body of work is that controllers conditioned on search state
yield more robust performance than those relying solely on cumulative operator
rewards.

\subsection{Reinforcement Learning for Combinatorial Optimization}

Reinforcement learning has become a major methodology for combinatorial
optimization. \citet{mazyavkina2021rl} survey RL methods across canonical
problems including TSP and vehicle routing. \citet{chung2025neural} extend this
taxonomy with an industrial-engineering perspective. On the construction-policy
side, \citet{khalil2017learning} combine graph embeddings with Deep Q-Networks
to learn greedy node-selection policies that generalize across instance sizes.
\citet{kool2019attention} propose an attention-based encoder--decoder achieving
near-optimal results for TSP instances up to 100 nodes. Unlike these neural
approaches, AdaptiSelect-TSP deliberately uses a compact tabular MDP, allowing
the learned policy and operator interactions to be inspected directly --- an
important property for interpretability and ablation.

% ─────────────────────────────────────────────────────────────────────────────
\section{Methodology}

\subsection{MDP Formulation}

We model operator selection as a discounted finite Markov Decision Process
$\langle \mathcal{S}, \mathcal{A}, \mathcal{R}, \mathcal{T}, \gamma \rangle$.
A state at iteration $t$ is the tuple $s_t = (q_t, \iota_t, \phi_t, a_{t-1})$
where:

\begin{itemize}
\item $q_t \in \{q_1,\ldots,q_5\}$ is the \textbf{solution quality band},
obtained by discretizing the gap to optimal into five bins: Excellent ($<$5\%),
Good (5--15\%), Fair (15--30\%), Poor (30--50\%), Very Poor ($>$50\%). When the
optimal is unknown, bands are defined relative to improvement from the initial
tour.

\item $\iota_t \in \{0, 1\}$ is the \textbf{recent improvement flag}, equal to
1 if the last step improved the tour and 0 otherwise.

\item $\phi_t \in \{\text{early}, \text{mid}, \text{late}\}$ is the
\textbf{iteration stage}, partitioning the budget $T$ into three equal thirds.
This feature enables true phase-awareness, allowing the policy to favour
exploration early and exploitation late.

\item $a_{t-1} \in \mathcal{A}$ is the \textbf{last operator used}, enabling
the controller to learn first-order operator-conditional preferences.
\end{itemize}

The Cartesian product yields $|\mathcal{S}| = 5 \times 2 \times 3 \times 4 =
120$ states. The action set is $\mathcal{A} = \{$2-opt, swap, relocate,
restart$\}$ and the discount factor is $\gamma = 0.9$.

\subsection{Local Search Operators}

The \textbf{2-opt} operator reverses a randomly chosen sub-segment of the tour;
a best-improvement variant samples $\min(2n, 100)$ candidate moves and retains
the best improving one. \textbf{Swap} exchanges two randomly selected cities.
\textbf{Relocate} removes a city and re-inserts it at a random position
(Or-opt style). \textbf{Restart} generates a completely random permutation.
An improving move is always accepted; a worsening move is accepted with
probability 0.01, inspired by Simulated Annealing, to prevent entrapment in
local optima.

\subsection{Reward Functions}

Two reward shapings are compared:
\begin{align}
R_1(s, a, s') &= L(\pi) - L(\pi') \label{eq:r1} \\
R_2(s, a, s') &= \bigl[L(\pi) - L(\pi')\bigr] + \beta \cdot \mathbf{1}[L(\pi') < L_{\text{best}}] \label{eq:r2}
\end{align}
where $L(\pi)$ is the tour length, $L_{\text{best}}$ is the running global best,
and $\beta = 100$ is the bonus magnitude. $R_1$ is the pure improvement reward;
$R_2$ adds a delayed signal that explicitly rewards global-best improvements.

\subsection{Q-Learning Update and Adaptive $\varepsilon$-Greedy}

The Q-table has dimensions $120 \times 4$, initialized to zero. Updates follow
the standard Bellman equation with learning rate $\alpha = 0.1$:
\begin{equation}
Q(s, a) \leftarrow Q(s, a) + \alpha \left[ r + \gamma \max_{a'} Q(s', a') - Q(s, a) \right]
\end{equation}

Three epsilon strategies are compared: (1)~\textbf{Constant} ($\varepsilon =
0.1$) as baseline; (2)~\textbf{Exponential Decay} starting at $\varepsilon =
1.0$ with decay rate 0.998; (3)~\textbf{Adaptive Stagnation} (proposed), where
$\varepsilon$ is multiplied by $\eta^+ = 1.5$ when consecutive non-improving
steps exceed threshold $\tau = 20$, and by $\eta^- = 0.9$ upon improvement.
This rule treats stagnation as evidence that the current exploitative policy is
locally insufficient and responds by injecting exploration.

\subsection{Operator Transition Matrix}

To quantify operator synergy we record an empirical matrix $M \in \mathbb{R}^{4
\times 4}$ where entry $M_{ij}$ estimates the conditional probability that
operator $a_j$ yields an improving move given that the previous successful
operator was $a_i$. Large $M_{ij}$ indicates that $a_i$ tends to enable
successful applications of $a_j$ --- i.e., the pair $(a_i, a_j)$ is
synergistic. The matrix is computed post-training as a diagnostic and is not
used inside the Q-update.

% ─────────────────────────────────────────────────────────────────────────────
\section{Experimental Setup}

We evaluate on four TSPLIB benchmark instances: eil51 ($n$=51, opt=426),
berlin52 ($n$=52, opt=7542), st70 ($n$=70, opt=675), and kroA100 ($n$=100,
opt=21282). All experiments use 5 independent random seeds, 3000 iterations per
episode, and results are reported as mean optimality gap (\%) $\pm$ standard
deviation. The initial tour is produced by the Nearest Neighbor (NN) heuristic.

Baselines include: Fixed 2-opt, Fixed Swap, Fixed Relocate, Fixed Restart (each
applying a single operator throughout), and Random (uniform selection at each
step). All methods share identical initialization, acceptance criterion, and
iteration budget. Hyperparameters: $\alpha$=0.1, $\gamma$=0.9, $\varepsilon_0$=0.3,
$\varepsilon_{\min}$=0.05, $\varepsilon_{\max}$=0.5, $\tau$=20, $\beta$=100.

% ─────────────────────────────────────────────────────────────────────────────
\section{Results and Discussion}

\subsection{Baseline Comparison}

Table~\ref{tab:main} reports the average optimality gap (\%) across 5 seeds for
all strategies and all four instances.

\begin{table}[h]
\centering
\caption{Average optimality gap (\%) $\pm$ std.\ dev.\ across 5 seeds. Lower is better. Best result per instance in \textbf{bold}.}
\label{tab:main}
\small
\begin{tabular}{lcccccc}
\toprule
Instance & \textbf{Q-Learning} & Fixed & Fixed & Fixed & Fixed & Random \\
         & \textbf{(Adaptive)} & 2-opt & Swap  & Relocate & Restart & \\
\midrule
eil51    & \textbf{3.00\,$\pm$\,1.07} & 3.94\,$\pm$\,1.55 & 22.35\,$\pm$\,3.19 & 20.56\,$\pm$\,2.21 & 24.04\,$\pm$\,4.25 & 3.99\,$\pm$\,0.71 \\
berlin52 & \textbf{1.03\,$\pm$\,1.01} & 2.17\,$\pm$\,2.41 & 17.12\,$\pm$\,4.17 & 16.39\,$\pm$\,4.00 & 17.50\,$\pm$\,4.30 & 1.28\,$\pm$\,2.67 \\
st70     & \textbf{3.82\,$\pm$\,2.06} & 5.42\,$\pm$\,1.80 & 27.70\,$\pm$\,5.31 & 25.48\,$\pm$\,6.47 & 27.97\,$\pm$\,5.51 & 6.99\,$\pm$\,2.69 \\
kroA100  & \textbf{3.90\,$\pm$\,0.76} & 5.10\,$\pm$\,0.62 & 27.72\,$\pm$\,3.59 & 28.18\,$\pm$\,4.13 & 28.27\,$\pm$\,4.07 & 16.68\,$\pm$\,3.13 \\
\bottomrule
\end{tabular}
\end{table}

On berlin52, Q-learning achieves a 1.03\% gap, compared to 2.17\% for fixed
2-opt --- a 52\% relative improvement. The advantage grows with instance size:
on kroA100, Q-learning achieves 3.90\% versus 5.10\% for fixed 2-opt and
16.68\% for random selection. Fixed swap, relocate, and restart perform very
poorly in isolation (17--28\% gaps), confirming that adaptive selection is
essential.

\subsection{Reward Shaping: R1 vs.\ R2}

\begin{table}[h]
\centering
\caption{Reward function comparison: average optimality gap (\%) $\pm$ std.\ dev., 5 seeds.}
\label{tab:reward}
\small
\begin{tabular}{lccc}
\toprule
Instance & R1 (Pure) & R2 (Global Bonus) & Fixed 2-opt \\
\midrule
eil51    & 3.12\,$\pm$\,1.20 & \textbf{2.95\,$\pm$\,0.98} & 3.94\,$\pm$\,1.55 \\
berlin52 & 1.18\,$\pm$\,1.30 & \textbf{1.03\,$\pm$\,1.01} & 2.17\,$\pm$\,2.41 \\
\bottomrule
\end{tabular}
\end{table}

R2 outperforms R1 on both instances and shows lower standard deviation (0.98
vs.\ 1.20 on eil51; 1.01 vs.\ 1.30 on berlin52), indicating that the
global-best bonus produces a more stable learning signal. Both reward functions
significantly outperform the fixed 2-opt baseline.

\subsection{Epsilon Strategy Comparison}

The adaptive stagnation-driven epsilon consistently achieves the best final gap
across all instances. Constant epsilon is competitive but shows slower early
convergence. Exponential decay converges faster initially due to high early
exploration, but lacks the ability to escape local optima in later stages.
The adaptive strategy's epsilon spikes upon stagnation detection are clearly
visible in training histories and consistently precede subsequent improvements.

\subsection{Operator Synergy Analysis}

The operator transition matrix on berlin52 reveals non-trivial emergent
strategies. The 2-opt $\rightarrow$ 2-opt self-loop dominates (97\% of
successful transitions from 2-opt), indicating self-reinforcing local
improvement chains. The Restart $\rightarrow$ 2-opt transition accounts for
100\% of post-restart successful moves: after a random restart, the agent
invariably applies 2-opt to rapidly improve the new tour. This recovery pattern
is not hard-coded --- it emerges entirely from Q-learning. Swap $\rightarrow$
2-opt (25\%) and Relocate $\rightarrow$ 2-opt (20\%) also show elevated rates,
suggesting these operators create configurations favourable to subsequent 2-opt
improvements.

\subsection{Generalization Test}

The Q-table trained on eil51 (10 episodes) was applied to st70 and kroA100 with
reduced learning rate ($\alpha$=0.05) and low epsilon ($\varepsilon$=0.05). The
pre-trained agent achieves gaps of 4.1\% on st70 and 4.3\% on kroA100, compared
to 3.82\% and 3.90\% for agents trained from scratch. The small gap
($<$0.5 percentage points) suggests the learned policy captures general search
dynamics rather than instance-specific patterns.

% ─────────────────────────────────────────────────────────────────────────────
\section{Conclusion}

We presented AdaptiSelect-TSP, a tabular Q-learning hyper-heuristic for
adaptive operator selection in TSP local search. The 120-state MDP explicitly
encodes search phase, enabling the policy to adapt operator preferences across
the trajectory without neural function approximation. Experimental results on
four TSPLIB instances demonstrate consistent outperformance of fixed and random
baselines, with the Q-learning agent achieving below 4\% optimality gap across
all tested instances. The global-best bonus reward (R2) and stagnation-driven
adaptive epsilon jointly account for the best performance. The operator
transition matrix reveals interpretable synergistic patterns --- particularly
the Restart $\rightarrow$ 2-opt recovery strategy --- that emerge autonomously
from Q-learning.

Future directions include: (i) Deep Q-Networks (DQN) for continuous state
representation and larger instances; (ii) incorporating Lin--Kernighan moves as
a fifth operator; (iii) extension to Capacitated VRP; and (iv) population-based
variants for maintaining solution diversity.

% ─────────────────────────────────────────────────────────────────────────────
\bibliography{references}
\bibliographystyle{iclr2026_conference}

\end{document}
